\scanpage{9}
où $\delta'$ \inserted{$\in\underline{\operatorname{NS}}_{B/S}(S)$} a été déjà considéré par ailleurs, par
\[
\widetilde{\delta'}\,\widetilde\delta=-\chi(\delta)\operatorname{id}_A\qquad\ldots..
\]
On peut se proposer, pour une section $\lambda$\\
de $\underline{\operatorname{Pic}}_{A/S}$, définissant $L$ faisceau inv. sur $A$,\\
de ``calculer'' la classe dans $K^\bullet(B)$ de\\
$\varphi_{\mathcal L}(\deleted{\unsure{\operatorname{cl}}}\operatorname{cl}(L))$ en termes de \deleted{\unsure{la même}} la section\\
$\varphi(\lambda)^{\unsure{\varepsilon|*}}$ de $\underline{\operatorname{Pic}}_{B/S}$, définissant un faisceau inversible\\
$L'^{-1}$ sur $B$ (de classe $-\delta'$ dans $\underline{\operatorname{NS}}(B/S)$ si\\
$\lambda$ est de classe $\delta$ dans $\operatorname{NS}(A/S)$). Heuristiquement,\\
on \deleted{aura} \inserted{devrait avoir} (pour $\chi(\delta)\neq0$ partout sur $S$)
\[
\text{``}\varphi_{\mathcal L}(\operatorname{cl}(L))=\chi\deleted{\unsure{\operatorname{cl}|\otimes}}L'^{-\frac1\chi}\text{''}
\qquad\text{où }\chi=\chi(\delta)\in H^0(S,\mathbf Z).
\]
\margin{indépendant\\
\unsure{des résultats|du résultat}\\
de Mumford}
Au niveau du $\operatorname{ch}_B$, ceci s'écrit
\[
\begin{aligned}
\operatorname{ch}\varphi_{\mathcal L}(\unsure{\operatorname{cl}L|\lambda})
&=\chi\exp\frac{-D'}\chi
=\chi-D'+\frac1{2\chi}D'^2-\frac1{3!\chi^2}D'^3\ldots\\
&\quad+(-1)^n\frac1{n!\chi^{n-\unsure{1}}}D'^n,
\end{aligned}
\]
\margin{$D'\in\operatorname{Gr}^1$\\
défini par $L'$}
et au niveau des classes de Chern \unsure{totales}
\[
\begin{aligned}
\operatorname{ct}(\varphi_{\mathcal L}(\unsure{\lambda|L}))&=\left(1-\frac{D'}\chi\right)^\chi\\
&=1-D'+\frac{\chi(\chi-1)}{2\chi}D'^2
-\frac{\chi(\chi-1)(\chi-2)}{3!\chi^2}D'^3+\ldots\\
&\quad+(-1)^n\frac{\chi(\chi-1)\ldots(\chi-n+1)}{n!\chi^{n-\unsure{1}}}D'^n.
\end{aligned}
\]
\margin{?}
ces formules sont-elles bien vraies mod torsion?

Elles sont vraies du moins sur un corps\\
alg clos, \uline{mod} \uuline{équivalence} \uline{cohomologique}\\
$l$-adique (tout $l\neq\operatorname{car}.k$). Ceci donne des\\
propriétés de divisibilité intéressantes de $\delta'$\\
\inserted{savoir} $\Rightarrow\deleted{\unsure{\chi}}\dfrac{(\chi-1)\ldots(\chi-k+1)}{k!\chi^{k-2}}\delta'^k$ est ``entier'' i.e. provient\\
d'un cycle à coeff. entiers
\margin{!\quad!!}
\mspage{8}
